\section{Additional Exercises}

\subsection{Vector Spaces}

\begin{exercise}\label{ex:multilinear-1}
Let $V:=\R^3$ be the set of all real triples.

We equip $V$ with addition $\oplus : V \times V\to V$ and s-multiplication $\odot:\R\times V \to V$ defined by
\bse
    (a,b,c)\oplus(d,e,f) := (a+d,b+e,c+f)
\ese
and
\bse
    \lambda \odot (a,b,c) := (\lambda \cdot a, \lambda \cdot b, \lambda \cdot c)
\ese
where $+$ and $\cdot$ are the addition and multiplication on $\R$. Check that $(V,\oplus,\odot)$ is a vector space.
\end{exercise}
\hyperref[sol:multilinear-1]{\small \textit{Solution on p.~\pageref*{sol:multilinear-1}}}

\begin{exercise}\label{ex:multilinear-2}
Consider the map $d:V\to V$; $(a,b,c)\mapsto d\big((a,b,c)\big) := (b,2c,0)$.

Is $d$ linear?
\end{exercise}
\hyperref[sol:multilinear-2]{\small \textit{Solution on p.~\pageref*{sol:multilinear-2}}}

\begin{exercise}\label{ex:multilinear-3}
Show that the map $d\circ d$ is linear.
\end{exercise}
\hyperref[sol:multilinear-3]{\small \textit{Solution on p.~\pageref*{sol:multilinear-3}}}

\begin{exercise}\label{ex:multilinear-4}
Consider the map
\bse
    i: V \to \R; \, (a,b,c) \mapsto i\big((a,b,c)\big) := a + \frac{1}{2}b + \frac{1}{3}c.
\ese
Check linearity. Of what set is $i$ an element?
\end{exercise}
\hyperref[sol:multilinear-4]{\small \textit{Solution on p.~\pageref*{sol:multilinear-4}}}

\begin{exercise}\label{ex:multilinear-5}
Another (multi)linear question. Again I don't want to keep typing out that check so see \href{https://www.youtube.com/watch?v=5oeWX3NUhMA&list=PLFeEvEPtX_0RQ1ys-7VIsKlBWz7RX-FaL&index=3}{the video} for this one.
\end{exercise}

\begin{exercise}\label{ex:multilinear-6}
Compare the above map $d:V \to V$ with the map $\del : P \to P$ from the lecture and construct a bijective linear map $j:P_2\to \R^3$ such that
\bse
    d = j \circ \del \circ j^{-1}.
\ese
\end{exercise}
\hyperref[sol:multilinear-6]{\small \textit{Solution on p.~\pageref*{sol:multilinear-6}}}

Similar exercises can be done to compare $i:V\to\R$ to $I:P\to\R$ given in the lecture (and also for the above exercise which I haven't typed here.) See \href{https://www.youtube.com/watch?v=5oeWX3NUhMA&list=PLFeEvEPtX_0RQ1ys-7VIsKlBWz7RX-FaL&index=3}{the video} for more details.

\subsection{Indices}

\begin{exercise}\label{ex:multilinear-7}
Let $V$ be a $d$-dimensional vector space. Consider two maps $A$ and $B$, where
\bse
    \begin{split}
        A : V^*\times V^* & \to \R \\
        B : V \times V & \to \R.
    \end{split}
\ese
$V$ has a basis $e_1,...,e_d$ and $V^*$ has the basis $\epsilon^1,...,\epsilon^d$.

Define the components $A^{ab}$ of $A$ and $B_{ab}$ of $B$ with respect to the given bases.
\end{exercise}
\hyperref[sol:multilinear-7]{\small \textit{Solution on p.~\pageref*{sol:multilinear-7}}}

\begin{exercise}\label{ex:multilinear-8}
We define $A^{[ab]} := \frac{1}{2}\big(A^{ab}-A^{ba}\big)$. Show that
\bse
    A^{[ab]} = - A^{[ba]}
\ese
and also
\bse
    A^{[ab]}B_{ab} = A^{ab}B_{[ab]}.
\ese
\end{exercise}
\hyperref[sol:multilinear-8]{\small \textit{Solution on p.~\pageref*{sol:multilinear-8}}}

\begin{exercise}\label{ex:multilinear-9}
We additionally define $B_{(ab)}:= \frac{1}{2}\big(B_{ab}+B_{ba}\big)$. Now, show that
\bse
    B_{(ab)} = B_{(ba)}
\ese
and again
\bse
    A^{ab}B_{(ab)} = A^{(ab)}B_{ab}.
\ese
\end{exercise}
\hyperref[sol:multilinear-9]{\small \textit{Solution on p.~\pageref*{sol:multilinear-9}}}

\begin{exercise}\label{ex:multilinear-10}
Using the results from the previous questions, we can easily show
\bse
    A^{[ab]}B_{(ab)} = 0,
\ese
i.e., the summation (contraction) of symmetric and antisymmetric indices yields zero.
\end{exercise}
\hyperref[sol:multilinear-10]{\small \textit{Solution on p.~\pageref*{sol:multilinear-10}}}

\subsection{Linear Maps As Tensors}

\begin{exercise}\label{ex:multilinear-11}
Given a vector space $V$ and a linear map $\phi: V^* \lmap V^*$ construct a $(1,1)$-tensor $T_{\phi}$.
\end{exercise}
\hyperref[sol:multilinear-11]{\small \textit{Solution on p.~\pageref*{sol:multilinear-11}}}

\begin{exercise}\label{ex:multilinear-12}
Given a $(1,1)$-tensor $T : V^*\times V \lmap \R$, construct a linear map $\phi_T :V^*\lmap V^*$.
\end{exercise}
\hyperref[sol:multilinear-12]{\small \textit{Solution on p.~\pageref*{sol:multilinear-12}}}

\begin{exercise}\label{ex:multilinear-13}
Show that
\ben[label=(\alph*)]
    \item $T_{\phi_T} = T$, and
    \item $\phi_{T_{\phi}}=\phi$.
\een
\end{exercise}
\hyperref[sol:multilinear-13]{\small \textit{Solution on p.~\pageref*{sol:multilinear-13}}}

\begin{exercise}\label{ex:multilinear-14}
Conclude that we can consider a linear map $\phi : V^*\lmap V^*$ as a $(1,1)$-tensor.
\end{exercise}
\hyperref[sol:multilinear-14]{\small \textit{Solution on p.~\pageref*{sol:multilinear-14}}}
